\documentclass[11pt]{article}
\usepackage{mathblatt}
% gesetzt von werkzeuge/setzer.py aus dem Lernweg (L5-A · L5-B · L5-C · L5-D · L5-T) und den Bankzeilen; Muster T6B.
% Präambel des Setzers (werkzeuge/setzer.py), wortgleich aus dem Hand-Satz T6B (bau/proben/2026-10-09/kathete-berechnen), dort aus K4W.
\makeatletter
\setlength{\mbpfbe}{0mm}% keine Punkte (Bauregel 6.4)
% eingerückt (Bauregel 4.7, 6.6): \gEin Einzug, \gSchrift eine Stufe kleiner
\newlength{\gEin}\setlength{\gEin}{0pt}
\let\gSchrift\relax
\renewcommand{\pfkreuzzeile}[1]{\par\noindent\hangindent3.2em\hangafter1\hspace*{1.6em}$\square$\ #1\par}
\newcommand{\gkreuz}[1]{$\square$\ #1\hspace{2em}}
% Bauregel 6.7: links, was man ansieht, rechts, was man tut; Spaltengrenze überall an derselben Stelle
\newsavebox{\gbox}\newlength{\gLinks}
\newcommand{\gzwei}[2]{\par\smallskip\noindent\sbox{\gbox}{#1}%
  \setlength{\gLinks}{\dimexpr0.5\textwidth-\gEin-\mbpfnr\relax}%
  \ifdim\wd\gbox>\dimexpr\gLinks-4mm\relax
    \usebox{\gbox}\par\smallskip\noindent\begin{minipage}{\linewidth}#2\end{minipage}%
  \else
    \begin{minipage}[t]{\gLinks}\vspace{0pt}\usebox{\gbox}\end{minipage}%
    \begin{minipage}[t]{\dimexpr\linewidth-\gLinks\relax}\vspace{0pt}\raggedright #2\end{minipage}%
  \fi\par}
\RenewDocumentEnvironment{pfaufg}{m m m}{%
  \begin{lrbox}{\mbaufgabebox}\begin{minipage}[t]{\linewidth}%
  \gSchrift\noindent\hspace*{\gEin}\mbpfrandmarke{#1}%
  \begin{minipage}[t]{\mbpfnr}\strut\textbf{#2}\end{minipage}%
  \begin{minipage}[t]{\dimexpr\linewidth-\gEin-\mbpfnr\relax}\raggedright\strut\ignorespaces}%
 {\end{minipage}%
  \end{minipage}\end{lrbox}%
  \setlength{\mbaufgabehoehe}{\dimexpr\ht\mbaufgabebox+\dp\mbaufgabebox\relax}%
  \par\addvspace{7pt}\Needspace*{\mbaufgabehoehe}%
  \noindent\usebox{\mbaufgabebox}\par\addvspace{7pt}}
\makeatother
% Rechenraum als Karo (Bauregel 6.9): n Rechenschritte = n mal 10 mm
\newcommand{\karo}[1]{\par\smallskip\noindent\begin{tikzpicture}
  \clip (0,0) rectangle ({\linewidth-1mm},{#1*10mm});
  \draw[black!16,line width=0.3pt,step=5mm] (0,0) grid ({\linewidth},{#1*10mm});
  \end{tikzpicture}\par}
\newcommand{\kk}[1]{$\square$\,#1\hspace{0.9em}}
\newcommand{\linie}[1]{\,{\color{black!45}\rule[-1pt]{#1}{0.4pt}}\,}
% Zeile am Ende jedes Durchgangs: Originale klein grau, darüber; dann Vorgänger/Nachfolger (Bauregel 3.4, 6.5)
\newcommand{\blattende}[2]{\par\vfill\noindent\begin{minipage}{\linewidth}{\scriptsize\color{mbgrau}Originale: #1\par}\smallskip
{\footnotesize #2\par}\end{minipage}}
\newcommand{\pkt}{\textperiodcentered{}}

\newcommand{\kennung}{HMF}
\newcommand{\grau}[1]{{\color{mbgrau}#1}}

\begin{document}
\pfheftstil
\fancyfoot[C]{\parbox[t]{\textwidth}{\mbfhbox\par\vspace{1.5mm}{\footnotesize\color{mbgrau}{\scriptsize \kennung}\hfill\thepage}}}
\pfheftkopf{Vergleichen, Ordnen, Runden}{Klasse 6}
% L5-A: Vergleichen und ordnen
\begin{pfaufg}{}{1.}{}
\grau{a)\ Ordne von klein nach groß: $0{,}5$; \ $0{,}45$; \ $0{,}405$; \ $0{,}54$.\par\smallskip Hänge Nullen an, bis alle drei Stellen nach dem Komma haben: $0{,}500$; \ $0{,}450$; \ $0{,}405$; \ $0{,}540$.\par Vergleiche von links. Einer: überall 0. Zehntel: 4 bei $0{,}450$ und $0{,}405$, 5 bei $0{,}500$ und $0{,}540$.\par Gleiche Zehntel: Die Hundertstel entscheiden. $0{,}405 < 0{,}450$ \ und \ $0{,}500 < 0{,}540$.\par Ohne die angehängten Nullen: $0{,}405 < 0{,}45 < 0{,}5 < 0{,}54$.}
\par\medskip
b)\ Setze $<$ oder $>$ ein.\par\vspace{3pt} a)~$0{,}58$ \leerfeld $0{,}54$ \qquad b)~$4{,}62$ \leerfeld $4{,}26$ \qquad c)~$15{,}2$ \leerfeld $16{,}2$ \qquad d)~$0{,}071$ \leerfeld $0{,}017$\karo{3}
\fusshilfe{\mbox{1:~b) a) $>$ \quad b) $>$ \quad c) $<$ \quad d) $>$}}
\end{pfaufg}

% L5-A: Vergleichen und ordnen
\begin{pfaufg}{}{2.}{}
Hänge Nullen an. Setze dann $<$, $=$ oder $>$ ein.\par\vspace{3pt} a)~$0{,}7$ \leerfeld $0{,}69$ \qquad b)~$3{,}45$ \leerfeld $3{,}5$ \qquad c)~$0{,}41$ \leerfeld $0{,}405$ \qquad d)~$5{,}1$ \leerfeld $5{,}10$\karo{3}
\fusshilfe{\mbox{2:~a) $>$ \quad b) $<$ \quad c) $>$ \quad d) $=$}}
\end{pfaufg}

% L5-A: Vergleichen und ordnen
\begin{pfaufg}{}{3.}{}
Ordne von klein nach groß: $4{,}3$; \ $4{,}03$; \ $4{,}33$; \ $4{,}303$.\karo{3}
\fusshilfe{\mbox{3:~$4{,}03 < 4{,}3 < 4{,}303 < 4{,}33$}}
\end{pfaufg}

% L5-A: Vergleichen und ordnen
\begin{pfaufg}{}{4.}{}
Ordne diesmal von groß nach klein: $5{,}08$; \ $5{,}8$; \ $5{,}008$; \ $5{,}088$.\karo{3}
\fusshilfe{\mbox{4:~$5{,}8 > 5{,}088 > 5{,}08 > 5{,}008$}}
\end{pfaufg}

% L5-A: Vergleichen und ordnen
\begin{pfaufg}{}{5.}{}
Fünf Kinder schwimmen $50$\,m. Die Zeiten: Ela $41{,}8$\,s, Finn $41{,}08$\,s, Gül $41{,}75$\,s, Hannes $42{,}1$\,s, Ida $41{,}1$\,s. Die schnellste Zeit gewinnt. Wer bekommt Silber (Platz 2), wer Bronze (Platz 3)?\karo{3}
\fusshilfe{\mbox{5:~Silber Ida ($41{,}1$\,s), Bronze Gül ($41{,}75$\,s)}}
\end{pfaufg}

% L5-B: Runden
\begin{pfaufg}{}{6.}{}
\grau{a)\ Runde $5{,}763$ auf Zehntel und auf Hundertstel. Runde $6{,}97$ auf Zehntel.\par\smallskip Auf Zehntel: Rechts vom Zehntel steht die 6. Bei 5 bis 9 aufrunden: $5{,}763 \approx 5{,}8$.\par Auf Hundertstel: Rechts vom Hundertstel steht die 3. Bei 0 bis 4 abrunden: $5{,}763 \approx 5{,}76$.\par $6{,}97$ auf Zehntel: Rechts steht die 7, also aufrunden: 9 Zehntel $+$ 1 Zehntel $=$ 10 Zehntel $=$ 1 Einer.\par Der Einer wird 7, an der Zehntelstelle bleibt die 0 stehen: $6{,}97 \approx 7{,}0$.}
\par\medskip
b)\ Runde auf Zehntel.\par\vspace{3pt} a)~$2{,}73$ \qquad b)~$5{,}48$ \qquad c)~$0{,}35$ \qquad d)~$8{,}04$\karo{3}
\fusshilfe{\mbox{6:~b) a) $2{,}7$ \quad b) $5{,}5$ \quad c) $0{,}4$ \quad d) $8{,}0$}}
\end{pfaufg}

% L5-B: Runden
\begin{pfaufg}{}{7.}{}
Runde auf Hundertstel.\par\vspace{3pt} a)~$1{,}234$ \qquad b)~$6{,}078$ \qquad c)~$12{,}465$ \qquad d)~$0{,}555$\karo{3}
\fusshilfe{\mbox{7:~a) $1{,}23$ \quad b) $6{,}08$ \quad c) $12{,}47$ \quad d) $0{,}56$}}
\end{pfaufg}

% L5-B: Runden
\begin{pfaufg}{}{8.}{}
Runde. Achte darauf, was aus 10 Zehnteln oder 10 Hundertsteln wird.\par\vspace{3pt} a)~$2{,}97$ auf Zehntel \par b)~$9{,}96$ auf Zehntel \par c)~$3{,}998$ auf Hundertstel \par d)~$0{,}95$ auf Zehntel\karo{3}
\fusshilfe{\mbox{8:~a) $3{,}0$ \quad b) $10{,}0$ \quad c) $4{,}00$ \quad d) $1{,}0$}}
\end{pfaufg}

% L5-B: Runden
\begin{pfaufg}{}{9.}{}
Runde die Größen.\par\vspace{3pt} a)~$2{,}346$\,€ auf Cent \par b)~$3{,}96$\,kg auf Zehntel \par c)~$1{,}784$\,m auf Zentimeter \par d)~$12{,}05$\,s auf Zehntel\karo{3}
\fusshilfe{\mbox{9:~a) $2{,}35$\,€ \quad b) $4{,}0$\,kg \quad c) $1{,}78$\,m \quad d) $12{,}1$\,s}}
\end{pfaufg}

% L5-B: Runden
\begin{pfaufg}{}{10.}{}
Mia kauft Käse. Die Waage rechnet einen Preis von $4{,}865$\,€ aus. Auf dem Bon steht der Preis auf Cent gerundet. Mia hat genau $4{,}85$\,€ dabei. Reicht ihr Geld?\karo{3}
\fusshilfe{\mbox{10:~Nein: Der Bon zeigt $4{,}87$\,€.}}
\end{pfaufg}

% L5-C: Dazwischen und genau in der Mitte
\begin{pfaufg}{}{11.}{}
\gzwei{\begin{tikzpicture}[font=\small]\draw[->,line width=0.7pt] (0,0) -- (6.65,0);\draw[line width=0.7pt] (0.000,-0.17) -- (0.000,0.17);\node[below] at (0.000,-0.17) {$2{,}3$};\draw[line width=0.4pt] (0.210,-0.1) -- (0.210,0.1);\draw[line width=0.4pt] (0.420,-0.1) -- (0.420,0.1);\draw[line width=0.4pt] (0.630,-0.1) -- (0.630,0.1);\draw[line width=0.4pt] (0.840,-0.1) -- (0.840,0.1);\draw[line width=0.4pt] (1.050,-0.13) -- (1.050,0.13);\draw[line width=0.4pt] (1.260,-0.1) -- (1.260,0.1);\draw[line width=0.4pt] (1.470,-0.1) -- (1.470,0.1);\draw[line width=0.4pt] (1.680,-0.1) -- (1.680,0.1);\draw[line width=0.4pt] (1.890,-0.1) -- (1.890,0.1);\draw[line width=0.7pt] (2.100,-0.17) -- (2.100,0.17);\node[below] at (2.100,-0.17) {$2{,}4$};\draw[line width=0.4pt] (2.310,-0.1) -- (2.310,0.1);\draw[line width=0.4pt] (2.520,-0.1) -- (2.520,0.1);\draw[line width=0.4pt] (2.730,-0.1) -- (2.730,0.1);\draw[line width=0.4pt] (2.940,-0.1) -- (2.940,0.1);\draw[line width=0.4pt] (3.150,-0.13) -- (3.150,0.13);\draw[line width=0.4pt] (3.360,-0.1) -- (3.360,0.1);\draw[line width=0.4pt] (3.570,-0.1) -- (3.570,0.1);\draw[line width=0.4pt] (3.780,-0.1) -- (3.780,0.1);\draw[line width=0.4pt] (3.990,-0.1) -- (3.990,0.1);\draw[line width=0.7pt] (4.200,-0.17) -- (4.200,0.17);\node[below] at (4.200,-0.17) {$2{,}5$};\draw[line width=0.4pt] (4.410,-0.1) -- (4.410,0.1);\draw[line width=0.4pt] (4.620,-0.1) -- (4.620,0.1);\draw[line width=0.4pt] (4.830,-0.1) -- (4.830,0.1);\draw[line width=0.4pt] (5.040,-0.1) -- (5.040,0.1);\draw[line width=0.4pt] (5.250,-0.13) -- (5.250,0.13);\draw[line width=0.4pt] (5.460,-0.1) -- (5.460,0.1);\draw[line width=0.4pt] (5.670,-0.1) -- (5.670,0.1);\draw[line width=0.4pt] (5.880,-0.1) -- (5.880,0.1);\draw[line width=0.4pt] (6.090,-0.1) -- (6.090,0.1);\draw[line width=0.7pt] (6.300,-0.17) -- (6.300,0.17);\node[below] at (6.300,-0.17) {$2{,}6$};\fill (3.150,0) circle (1.7pt);\node[above] at (3.150,0.12) {$2{,}45$};\end{tikzpicture}}{\grau{a)\ Welche Zahl liegt genau in der Mitte zwischen $2{,}3$ und $2{,}6$? Nenne noch eine Zahl, die dazwischen liegt.\par\smallskip Hänge an beide Zahlen eine Null an: $2{,}30$ und $2{,}60$.\par Lies sie ohne Komma (in Hundertsteln): $230$ und $260$.\par Abstand: $260 - 230 = 30$. Die Hälfte: $15$. Mitte: $230 + 15 = 245$.\par Komma zurück, so viele Stellen wie nach dem Anhängen: Mitte $= 2{,}45$.\par Jede Zahl zwischen $230$ und $260$ liegt dazwischen, z.\,B. $231$, also $2{,}31$.}}
\par\medskip
b)\ Gib eine Zahl an, die dazwischen liegt.\par\vspace{3pt} a)~$0{,}3$ und $0{,}4$ \qquad b)~$3{,}4$ und $3{,}5$ \qquad c)~$0{,}25$ und $0{,}26$\karo{3}
\fusshilfe{\mbox{11:~b) a) z.\,B}}
\end{pfaufg}

% L5-C: Dazwischen und genau in der Mitte
\begin{pfaufg}{}{12.}{}
Welche Zahl liegt genau in der Mitte?\par\vspace{3pt} a)~$0{,}6$ und $0{,}7$ \qquad b)~$4{,}1$ und $4{,}2$ \qquad c)~$0{,}04$ und $0{,}05$\karo{3}
\fusshilfe{\mbox{12:~a) $0{,}65$ \quad b) $4{,}15$ \quad c) $0{,}045$}}
\end{pfaufg}

% L5-C: Dazwischen und genau in der Mitte
\begin{pfaufg}{}{13.}{}
\gzwei{\begin{tikzpicture}[font=\small]\draw[->,line width=0.7pt] (0,0) -- (6.35,0);\draw[line width=0.7pt] (0.000,-0.17) -- (0.000,0.17);\node[below] at (0.000,-0.17) {$1$};\draw[line width=0.4pt] (0.600,-0.1) -- (0.600,0.1);\draw[line width=0.4pt] (1.200,-0.1) -- (1.200,0.1);\draw[line width=0.4pt] (1.800,-0.1) -- (1.800,0.1);\draw[line width=0.4pt] (2.400,-0.1) -- (2.400,0.1);\draw[line width=0.4pt] (3.000,-0.13) -- (3.000,0.13);\draw[line width=0.4pt] (3.600,-0.1) -- (3.600,0.1);\draw[line width=0.4pt] (4.200,-0.1) -- (4.200,0.1);\draw[line width=0.4pt] (4.800,-0.1) -- (4.800,0.1);\draw[line width=0.4pt] (5.400,-0.1) -- (5.400,0.1);\draw[line width=0.7pt] (6.000,-0.17) -- (6.000,0.17);\node[below] at (6.000,-0.17) {$2$};\fill (1.200,0) circle (1.7pt);\node[above] at (1.200,0.12) {$1{,}2$};\fill (4.800,0) circle (1.7pt);\node[above] at (4.800,0.12) {$1{,}8$};\end{tikzpicture}}{%
Markiere am Zahlenstrahl und schreibe die Zahl auf.\par\vspace{3pt} a)~die Mitte zwischen $1{,}2$ und $1{,}8$ \par b)~die Mitte zwischen $1{,}2$ und $1{,}5$\karo{3}}
\fusshilfe{\mbox{13:~a) $1{,}5$ \quad b) $1{,}35$}}
\end{pfaufg}

% L5-C: Dazwischen und genau in der Mitte
\begin{pfaufg}{}{14.}{}
Welche Zahl liegt genau in der Mitte? Hier liegt ein Einer dazwischen.\par\vspace{3pt} a)~$2{,}9$ und $3{,}2$ \qquad b)~$5{,}7$ und $6{,}4$\karo{3}
\fusshilfe{\mbox{14:~a) $3{,}05$ \quad b) $6{,}05$}}
\end{pfaufg}

% L5-C: Dazwischen und genau in der Mitte
\begin{pfaufg}{}{15.}{}
An einem Wanderweg steht eine Hütte bei km $3{,}4$ und eine Brücke bei km $4{,}1$. Genau in der Mitte dazwischen soll eine Bank stehen. Ein Rastplatz liegt bei km $3{,}8$. Liegt der Rastplatz zwischen Hütte und Bank oder zwischen Bank und Brücke?\karo{3}
\fusshilfe{\mbox{15:~Bank bei km $3{,}75$}}
\end{pfaufg}

% L5-D: Bruch, Prozent, Quadrat: erst umwandeln
\begin{pfaufg}{}{16.}{}
\grau{a)\ Welcher Wert ist der kleinste: $\dfrac{3}{20}$; \ $0{,}3^2$; \ $5\,\%$; \ $0{,}3$?\par\smallskip Bruch: Erweitere auf 100: $\dfrac{3}{20} = \dfrac{15}{100} = 0{,}15$.\par Quadrat: $0{,}3^2 = 0{,}3 \cdot 0{,}3$. Rechne $3 \cdot 3 = 9$. Beide Zahlen haben zusammen 2 Stellen nach dem Komma: $0{,}09$.\par Prozent heißt Hundertstel: $5\,\% = \dfrac{5}{100} = 0{,}05$.\par Vergleiche wie in A: $0{,}15$; \ $0{,}09$; \ $0{,}05$; \ $0{,}30$. Am kleinsten: $0{,}05$, also $5\,\%$.}
\par\medskip
b)\ Schreibe als Dezimalzahl.\par\vspace{3pt} a)~$\dfrac{3}{5}$ \qquad b)~$45\,\%$ \qquad c)~$0{,}2^2$ \qquad d)~$7\,\%$ \qquad e)~$\dfrac{3}{4}$ \qquad f)~$0{,}6^2$\karo{3}
\fusshilfe{\mbox{16:~b) a) $0{,}6$ \quad b) $0{,}45$ \quad c) $0{,}04$ \quad d) $0{,}07$ \quad e) $0{,}75$ \quad f) $0{,}36$}}
\end{pfaufg}

% L5-D: Bruch, Prozent, Quadrat: erst umwandeln
\begin{pfaufg}{}{17.}{}
Setze $<$, $=$ oder $>$ ein.\par\vspace{3pt} a)~$\dfrac{3}{4}$ \leerfeld $0{,}7$ \qquad b)~$8\,\%$ \leerfeld $0{,}8$ \qquad c)~$0{,}7^2$ \leerfeld $0{,}7$ \qquad d)~$\dfrac{3}{20}$ \leerfeld $15\,\%$\karo{3}
\fusshilfe{\mbox{17:~a) $>$ \quad b) $<$ \quad c) $<$ \quad d) $=$}}
\end{pfaufg}

% L5-D: Bruch, Prozent, Quadrat: erst umwandeln
\begin{pfaufg}{}{18.}{}
Ordne von klein nach groß: $\dfrac{2}{5}$; \ $0{,}38$; \ $39\,\%$; \ $0{,}6^2$.\karo{3}
\fusshilfe{\mbox{18:~$0{,}6^2 < 0{,}38 < 39\,\% < \dfrac{2}{5}$}}
\end{pfaufg}

% L5-D: Bruch, Prozent, Quadrat: erst umwandeln
\begin{pfaufg}{}{19.}{}
Kreuze die Aussage an, die wahr ist.\par\medskip\noindent\hspace*{7mm}\mbox{$\square$\ $0{,}3^2 = 0{,}6$}\hspace{10mm}\mbox{$\square$\ $\dfrac{1}{4} > 0{,}3$}\hspace{10mm}\mbox{$\square$\ $6\,\% = 0{,}06$}\hspace{10mm}\mbox{$\square$\ $0{,}9 < 85\,\%$}\rule[-3mm]{0pt}{8mm}\par\smallskip 
\fusshilfe{\mbox{19:~$6\,\% = 0{,}06$}}
\end{pfaufg}

% L5-D: Bruch, Prozent, Quadrat: erst umwandeln
\begin{pfaufg}{}{20.}{}
Drei Klassen zählen, wer mit dem Rad zur Schule kommt. In der 6a sind es $\dfrac{3}{4}$ der Kinder, in der 6b $0{,}68$ der Kinder, in der 6c $72\,\%$ der Kinder. In welcher Klasse ist der Anteil am größten?\par\smallskip \gkreuz{6a}\gkreuz{6b}\par \gkreuz{6c}
\fusshilfe{\mbox{20:~6a ($0{,}75$)}}
\end{pfaufg}

% L5-T: Probetest
\begin{pfaufg}{}{21.}{}
Ordne von klein nach groß: $2{,}08$; \ $2{,}8$; \ $2{,}078$; \ $2{,}18$.\karo{3}
\fusshilfe{\mbox{21:~$2{,}078 < 2{,}08 < 2{,}18 < 2{,}8$}}
\end{pfaufg}

% L5-T: Probetest
\begin{pfaufg}{}{22.}{}
Runde.\par\vspace{3pt} a)~$7{,}348$ auf Zehntel \qquad b)~$7{,}348$ auf Hundertstel \qquad c)~$9{,}97$ auf Zehntel\karo{3}
\fusshilfe{\mbox{22:~a) $7{,}3$ \quad b) $7{,}35$ \quad c) $10{,}0$}}
\end{pfaufg}

% L5-T: Probetest
\begin{pfaufg}{}{23.}{}
Welche Zahl liegt genau in der Mitte zwischen $4{,}7$ und $4{,}8$?\karo{3}
\fusshilfe{\mbox{23:~$4{,}75$}}
\end{pfaufg}

% L5-T: Probetest
\begin{pfaufg}{}{24.}{}
Setze $<$, $=$ oder $>$ ein.\par\vspace{3pt} a)~$\dfrac{7}{20}$ \leerfeld $0{,}35$ \qquad b)~$6\,\%$ \leerfeld $0{,}6$ \qquad c)~$0{,}9^2$ \leerfeld $0{,}9$\karo{3}
\fusshilfe{\mbox{24:~a) $=$ \quad b) $<$ \quad c) $<$}}
\end{pfaufg}

% L5-T: Probetest
\begin{pfaufg}{}{25.}{}
Welcher Wert ist der kleinste? Kreuze an.\par\smallskip \gkreuz{$0{,}8$}\gkreuz{$0{,}08$}\par \gkreuz{$0{,}2^2$}\gkreuz{$7\,\%$}
\fusshilfe{\mbox{25:~$0{,}2^2 = 0{,}04$}}
\end{pfaufg}

% L5-T: Probetest
\begin{pfaufg}{}{26.}{}
Für das Klassenfest gibt es drei Krüge mit Saft und Sirup. In Krug A ist $\dfrac{1}{4}$ des Getränks Sirup, in Krug B $0{,}2$ und in Krug C $22\,\%$. Lea mag es möglichst wenig süß. Welchen Krug nimmt sie? Welcher Krug ist am süßesten?\karo{3}
\fusshilfe{\mbox{26:~Lea}}
\end{pfaufg}

\par\vfill\noindent{\footnotesize Vorher: Dezimalzahlen\quad\textbullet\quad Weiter: Rechnen mit Brüchen (Bruchrechnung)\par}
\end{document}
